% Part: sets-functions-relations
% Chapter: sets
% Section: proofs-about-sets

\documentclass[../../../include/open-logic-section]{subfiles}

\begin{document}

\olfileid{sfr}{set}{prf}
\olsection{برهان‌هایی دربارهٔ مجموعه‌ها}

\begin{editorial}
این بخش با~\verb|content/method/proofs| جایگزین شده و در حالت پیش‌فرض از
این فصل حذف شده است.
\end{editorial}

\begin{explain}
مجموعه‌ها و نمادگذاری‌هایی که تاکنون معرفی کرده‌ایم، صورت‌های کوتاه و
مناسبی برای مشخص‌کردن مجموعه‌ها و بیان رابطه‌های میان آن‌ها فراهم می‌کنند.
اغلب لازم است ادعاهایی دربارهٔ این رابطه‌ها را نیز اثبات کنیم. اگر با
برهان‌های ریاضی آشنا نباشید، این کار ممکن است برایتان تازه باشد. ازاین‌رو
نمونه‌ای ساده را گام‌به‌گام بررسی می‌کنیم. ثابت می‌کنیم برای هر دو مجموعهٔ
$X$ و~$Y$ همواره~$X \cap (X \cup Y) = X$. چگونه همانی‌ای میان مجموعه‌ها
مانند این را اثبات می‌کنیم؟ به یاد آورید که مجموعه‌ها تنها با~!!{element}s
خود تعیین می‌شوند؛ یعنی دو مجموعه همان‌اند اگر و تنها اگر~!!{element}s
یکسانی داشته باشند. پس در این مورد باید نشان دهیم (الف) هر~!!{element}
از~$X \cap (X \cup Y)$، !!{element}ی از~$X$ نیز هست و، برعکس، (ب)
هر~!!{element} از~$X$، !!{element}ی از~$X \cap (X \cup Y)$ نیز هست.
به بیان دیگر، هم (الف)~$X \cap (X \cup Y) \subseteq X$ و هم
(ب)~$X \subseteq X \cap (X \cup Y)$ را نشان می‌دهیم.

برهان ادعایی کلی مانند «هر~!!{element}~$z$ از~$X \cap (X \cup Y)$،
!!{element}ی از~$X$ نیز هست» با این کار آغاز می‌شود که~$z \in
X \cap (X \cup Y)$ دلخواهی را مفروض بگیریم و از این فرض نتیجه بگیریم
$z \in X$. شاید این الگو را با نام «برهان شرطی کلی» بشناسید. در این
برهان باید تعریف‌های دخیل در فرض و نتیجه، در اینجا تعریف‌های~«$\cap$» و
«$\cup$»، را نیز به کار ببریم. پس مورد~(الف) چنین استدلال می‌شود:

\begin{quote}
(الف) نخست می‌خواهیم نشان دهیم~$X \cap (X \cup Y) \subseteq X$؛ یعنی
بنا بر تعریف~$\subseteq$، برای هر~$z$، اگر~$z \in X \cap (X \cup Y)$،
آنگاه~$z \in X$. پس فرض کنید~$z \in X \cap (X \cup Y)$. چون~$z$
!!{element}ی از اشتراک دو مجموعه است اگر و تنها اگر~!!{element}ی از هر
دو باشد، نتیجه می‌گیریم~$z \in X$ و نیز~$z \in X \cup Y$. به‌ویژه،
$z \in X$، و این همان چیزی است که می‌خواستیم نشان دهیم.
\end{quote}

این نیمهٔ نخست برهان را کامل می‌کند. توجه کنید در گام آخر از این واقعیت
استفاده کردیم که اگر عطفی (~$z \in X$ و~$z \in X \cup Y$) از یک فرض
نتیجه شود، هر مؤلفهٔ عطفی نیز از همان فرض نتیجه می‌شود. شاید این قاعده را
با نام «حذف عطف» یا~$\land$Elim بشناسید. اکنون~(ب) را اثبات می‌کنیم:

\begin{quote}
(ب) اکنون~$X \subseteq X \cap (X \cup Y)$ را اثبات می‌کنیم؛ یعنی بنا
بر تعریف~$\subseteq$، برای هر~$z$، اگر~$z \in X$، آنگاه
$z \in X \cap (X \cup Y)$. فرض کنید~$z \in X$. برای نشان‌دادن
$z \in X \cap (X \cup Y)$ باید بنا بر تعریف~«$\cap$» نشان دهیم
(۱)~$z \in X$ و نیز (۲)~$z \in X \cup Y$. بند~(۱) همان فرض ماست، پس
چیز دیگری برای اثبات ندارد. برای~(۲)، به یاد آورید که~$z$ !!{element}ی
از اجتماع مجموعه‌هاست اگر و تنها اگر~!!{element}ی از دست‌کم یکی از آن
مجموعه‌ها باشد. چون~$z \in X$ و~$X \cup Y$ اجتماع~$X$ و~$Y$ است، این
شرط در اینجا برقرار است. پس~$z \in X \cup Y$. هم (۱)~$z \in X$ و هم
(۲)~$z \in X \cup Y$ را نشان داده‌ایم؛ بنابراین بنا بر تعریف~«$\cap$»،
$z \in X \cap (X \cup Y)$.
\end{quote}

این استدلال قدری طولانی بود، اما چگونگی استدلال دربارهٔ مجموعه‌ها و
رابطه‌های میان آن‌ها را نشان می‌دهد. معمولاً تا این اندازه تصریح نمی‌کنیم؛
به‌ویژه، شاید همهٔ تعریف‌ها را تکرار نکنیم. برهان این نتیجه در متنی
پیشرفته‌تر بسیار فشرده‌تر است و ممکن است چنین باشد.
\end{explain}


\begin{prop}[جذب]
برای همهٔ مجموعه‌های~$X$ و~$Y$،
\[
X \cap (X \cup Y) = X
\]
\end{prop}

\begin{proof}
(الف) فرض کنید~$z \in X \cap (X \cup Y)$. آنگاه~$z \in X$؛ پس
$X \cap (X \cup Y) \subseteq X$.

(ب) اکنون فرض کنید~$z \in X$. آنگاه~$z \in X \cup Y$ و در نتیجه
$z \in X \cap (X \cup Y)$. پس~$X \subseteq X \cap (X \cup Y)$.
\end{proof}

\begin{prob}
با جزئیات ثابت کنید~$X \cup (X \cap Y) = X$. سپس برهانی کوتاه و فشرده
ارائه کنید. (راهنمایی: برای جهت~$X \cup (X \cap Y) \subseteq X$ به
برهان بر حسب حالت‌ها، یا~$\lor$Elim، نیاز دارید.)
\end{prob}

\end{document}
